\documentclass[aspectratio=169,11pt]{ctexbeamer}

\usetheme{Madrid}
\usecolortheme{default}
\usefonttheme{professionalfonts}
\setbeamertemplate{navigation symbols}{}

\usepackage{amsmath,amssymb,amsthm,mathrsfs,bm}
\usepackage{tikz}
\usepackage{graphicx}
\usepackage{booktabs}
\usepackage{mathtools}
\usepackage{hyperref}

\newcommand{\A}{\mathbb{A}}
\newcommand{\C}{\mathbb{C}}
\newcommand{\D}{\mathcal{D}}
\newcommand{\OO}{\mathcal{O}}
\newcommand{\gr}{\operatorname{gr}}
\newcommand{\Ann}{\operatorname{Ann}}
\newcommand{\Hom}{\operatorname{Hom}}
\newcommand{\Supp}{\operatorname{Supp}}
\newcommand{\Char}{\operatorname{Char}}

\AtBeginSection[]{
    \begin{frame}[plain]
        \centering
        \vfill
        {\Large 第 \insertsectionnumber\ 部分\par}
        \vspace{1em}
        {\huge\bfseries \insertsection\par}
        \vfill
    \end{frame}
}


\title{D模理论综述}

\author{姚远}
\institute[海南大学数学与统计学院]{海南大学数学与统计学院\\[0.6em]指导老师：陈颖}



\date{2026年5月23日}

\begin{document}






%------------------------------------------------
\begin{frame}
    \titlepage
\end{frame}

%------------------------------------------------
\begin{frame}{报告结构}
    \tableofcontents
\end{frame}

\section{研究背景与理论框架}

%------------------------------------------------
\begin{frame}{论文研究内容}
    本文在特征零域上的仿射情形中，研究 Weyl 代数 $A_n$ 及其左模的结构性质。

    \vspace{0.6em}
    主要内容包括：

    \begin{enumerate}
        \item 从多项式环上的微分算子构造 Weyl 代数 $A_n$；
        \item 利用 Bernstein 滤过和良好滤过定义 $A_n$-模的维数与重数；
        \item 证明 Bernstein 不等式 $d(M)\ge n$，并讨论 holonomic 模的有限性；
        \item 介绍 Bernstein--Sato 多项式的存在性、计算方法及其在奇点理论中的作用。
    \end{enumerate}
    \pause
    \begin{block}{研究目标}
        在纯代数框架下刻画 $A_n$-模的增长行为与结构复杂性，并说明其与 D-模理论和奇点分析的联系。
    \end{block}

\end{frame}

%------------------------------------------------
\begin{frame}{从微分方程到代数对象}
    多项式系数的线性（偏）微分方程可以写成
    \[
        P(x,\partial)u=0.
    \]
    \pause
    例如
    \[
        \frac{du}{dx}=xu
        \quad\Longleftrightarrow\quad
        (\partial_x-x)u=0.
    \]

    \vspace{0.6em}
    \pause
    因而，微分方程可以转化为微分算子环上的代数问题。

    \[
        \text{微分方程}
        \quad\Longrightarrow\quad
        \text{微分算子的代数}
    \]

\end{frame}


%------------------------------------------------
\begin{frame}{D-模：微分方程的代数化}
    设 $\D_X$ 是空间 $X$ 上的微分算子层。
    \pause
    \begin{block}{定义}
        D-模就是 $\D_X$-模；也就是允许微分算子作用的对象。
    \end{block}
    \pause
    一个方程
    \[
        P u=0
    \]
    可以对应为
    \[
        M=\D_X/\D_XP.
    \]
    \pause
    \begin{block}{直观理解}
        D-模不直接给出方程的显式解，而是记录解所满足的微分关系，从而便于研究解空间结构及其几何性质。
    \end{block}

    \vspace{0.6em}
    因此，即使不能显式求解，也可以讨论解空间、奇点、延拓和单值性等结构问题。

\end{frame}

%------------------------------------------------
\begin{frame}{滤过、维数与 holonomic 模}
    为了衡量 D-模的复杂度，引入 Bernstein 滤过：
    \[
        B_m=\operatorname{span}\{x^\alpha\partial^\beta:|\alpha|+|\beta|\le m\}.
    \]

    \begin{itemize}
        \item 滤过：按“算子复杂度”分层；
        \item Hilbert 多项式：描述复杂度随阶数如何增长；
        \item 维数 $d(M)$：增长速度的代数刻画。
    \end{itemize}
    \pause
    \begin{block}{Holonomic 模}
        Bernstein 不等式给出 $d(M)\ge n$。当 $d(M)=n$ 时，称为 holonomic 模。
    \end{block}
    \pause
    \begin{block}{直观理解}
        Holonomic 模达到 Bernstein 不等式的下界，是增长速度最小、结构最有限的一类 D-模。
    \end{block}

\end{frame}



%------------------------------------------------
\begin{frame}{从奇点问题出发}
    设：
    \[
        f(x)\in\C[x_1,\dots,x_n].
    \]

    考虑多值函数：
    \[
        f^s.
    \]

    其中 $s$ 为复参数。研究 $f^s$ 的微分方程关系，是连接 D-模与奇点理论的重要入口。

    问题：

    \begin{itemize}
        \item 如何研究 $f^s$ 在奇点附近的行为？
        \item 如何描述其解析延拓？
        \item 奇点如何影响微分方程？
    \end{itemize}

    \begin{block}{核心问题}
        如何用代数方法刻画 $f=0$ 附近的奇点、解析延拓和单值性？
    \end{block}

\end{frame}



\section{Bernstein--Sato 多项式}

%------------------------------------------------
\begin{frame}{Bernstein--Sato 方程}
    Bernstein 的基本结论是：$f^s$ 满足一个由微分算子控制的函数方程。
    \pause
    \begin{block}{Bernstein--Sato 方程}
        存在微分算子 $P(s)$ 与多项式 $b(s)$ 使得
        \[
            P(s)f^{s+1}=b(s)f^s.
        \]
    \end{block}

    满足该关系的首一最小多项式称为 Bernstein--Sato 多项式，也称 $b$-函数。
    \pause
    $b_f(s)$ 是由微分算子作用关系提取出的奇点不变量；它将局部几何信息压缩到一元多项式中。


\end{frame}

%------------------------------------------------
\begin{frame}{b-函数的基本性质}
    Bernstein--Sato 多项式具有以下基本性质：

    \begin{enumerate}
        \item $b(s)$ 的根都是负有理数；
        \item $-1$ 总是一个根；
        \item 根与奇点单值性相关；
        \item 与 Hodge 理论、奇点理论和局部上同调密切相关。
    \end{enumerate}

    \begin{block}{意义}
        这些性质说明 $b_f(s)$ 不是形式上的辅助多项式，而是反映奇点结构的代数不变量。
    \end{block}

\end{frame}

%------------------------------------------------
\begin{frame}{b-函数的基本例子}
    当 $f(x)=x$ 时，
    \[
        \partial_x x^{s+1}=(s+1)x^s,
        \qquad b_f(s)=s+1.
    \]
    \pause
    \vspace{0.6em}
    当 $f(x)=x^m$ 时，
    \[
        b_f(s)=\prod_{k=1}^m\left(s+\frac{k}{m}\right).
    \]
    \pause
    \begin{block}{观察}
        从 $x$ 到 $x^m$，根的数量和分母发生变化，反映了零点多重性和局部分支结构的增加。
    \end{block}

\end{frame}

%------------------------------------------------
\begin{frame}{b-函数作为奇点不变量}
    $b_f(s)$ 的根不是随机出现的数，它们反映了 $f=0$ 的局部几何。

    \begin{itemize}
        \item 根的位置：反映奇点的复杂程度；
        \item 根的分母：常与分支、覆盖、单值性有关；
        \item 最大根：与 log canonical threshold 等不变量相关；
        \item 重根现象：提示奇点附近可能有更复杂的层次结构。
    \end{itemize}

    \vspace{0.8em}
    \begin{block}{概括}
        $b_f(s)$ 通过微分方程和 D-模结构刻画 $f=0$ 附近的局部几何，是奇点理论中的重要代数不变量。
    \end{block}

\end{frame}



\section{应用与进一步讨论}



%------------------------------------------------
\begin{frame}{应用一：区分不同类型的奇点}
    除了光滑点和多重点，b-函数还能识别更细的奇点结构。
    \pause
    \begin{columns}[T]
        \begin{column}{0.48\textwidth}
            节点（node） $f=x^2+y^2$：
            \[
                b_f(s)=(s+1)^2.
            \]
            这是两个光滑分支横截相交的情形。
        \end{column}
        \pause
        \begin{column}{0.48\textwidth}
            尖点奇点 （cusp）$f=x^2+y^3$：
            \[
                b_f(s)=(s+1)\left(s+\frac56\right)\left(s+\frac76\right).
            \]
            出现非整数根，反映尖点更复杂的单值性与局部拓扑信息。
        \end{column}
    \end{columns}
    \pause
    \vspace{0.8em}
    \begin{block}{说明}
        节点和尖点在局部几何上不同；$b$-函数通过根的差异记录这种局部结构差异。
    \end{block}

\end{frame}





%------------------------------------------------
\begin{frame}{应用二：单值性与 Milnor 纤维}
    考虑 Milnor 纤维：
    \[
        F_\varepsilon=f^{-1}(\varepsilon).
    \]

    围绕奇点作解析延拓，会得到单值变换：
    \[
        T:H^*(F_\varepsilon)\to H^*(F_\varepsilon).
    \]
    \pause
    Malgrange--Kashiwara 定理给出如下联系：

    \begin{block}{重要联系}
        若 $\alpha$ 是 $b_f(s)$ 的根，则
        \[
            e^{-2\pi i\alpha}
        \]
        是单值变换的特征值。
    \end{block}

    因此：

    \begin{center}
        \Large
        b-函数 $\Longleftrightarrow$ 奇点拓扑
    \end{center}

    这表明 $b$-函数不仅包含代数信息，也反映奇点附近的拓扑行为。
\end{frame}

\begin{frame}{一个具体单值性例子}
    取
    \[
        f(x)=x^m.
    \]
    则
    \[
        b_f(s)=\prod_{k=1}^m\left(s+\frac{k}{m}\right).
    \]

    纤维
    \[
        f^{-1}(t)=\{t^{1/m}e^{2\pi i k/m}:k=0,\dots,m-1\}
    \]
    由 $m$ 个点组成。

    \begin{block}{单值性现象}
        当 $t$ 绕原点一圈，这 $m$ 个点循环置换；对应特征值为 $m$ 次单位根，与 $b$-函数根给出的指数信息相匹配。
    \end{block}

\end{frame}
%------------------------------------------------






\section{总结}

%------------------------------------------------
\begin{frame}{总结}
    本文主线：

    \begin{center}
        微分方程
        $\Longrightarrow$
        微分算子环
        $\Longrightarrow$
        Weyl 代数
        $\Longrightarrow$
        Holonomic 模
        $\Longrightarrow$
        b-函数与奇点理论
    \end{center}

    \vspace{1em}
    主要结论：

    \begin{itemize}
        \item Weyl 代数为线性微分方程提供了非交换代数模型；
        \item Bernstein 滤过与 Hilbert 多项式刻画了 $A_n$-模的增长复杂度；
        \item Bernstein 不等式与 holonomic 模揭示了 D-模中的有限性现象；
        \item $b$-函数连接了 D-模结构、奇点理论、单值性和算法计算。
    \end{itemize}

\end{frame}

%------------------------------------------------
\begin{frame}{进一步研究方向}
    基于本文内容，可以继续从以下方向展开：

    \begin{itemize}
        \item 将仿射情形推广到一般光滑代数簇或复流形；
        \item 深入研究 Riemann--Hilbert 对应与 perverse sheaves；
        \item 结合 Hodge 模理论研究更精细的奇点不变量。
    \end{itemize}

\end{frame}





%------------------------------------------------
\begin{frame}
    \centering
    \Huge 谢谢！

    \vspace{1em}
    \Large Questions?
\end{frame}

\end{document}
